% Cross-model robustness comparison (V6).
% Reads the two focus phases vertically by model. V6 deltas: margins restated
% from tab_b1_gap (the V5 "under 1.5 pp from L3 onward" claim was wrong for
% the CNN-mean comparison: the L3 gap is +2.05 pp vs the mean and +0.94 vs
% the better CNN); MNIST explicitly stated as no-advantage territory; plain
% language reading woven in.
\subsection{Cross-model robustness comparison}
\label{sec:crossmodel}

This subsection reads the two focus phases vertically, by model. The
combined-degradation behaviour comes from Phase~B1
(Tables~\ref{tab:phase_b}--\ref{tab:b1_gap}, Figure~\ref{fig:b1_curves}), the
single-axis behaviour from Phase~C2 (Table~\ref{tab:phase_c2},
Figures~\ref{fig:c2_type_psnr}--\ref{fig:c2_model_ssim}).

\subsubsection{Per-model profiles}
\label{sec:cm:profiles}

\sectheading{ResNet50 --- the fast, ordinary baseline.} ResNet50 opens at
95.18\% clean on CIFAR-10 and 99.14\% on MNIST, and declines monotonically
under the combined recipe to 19.06\% at L5 on CIFAR-10. Its Phase~C2 profile
is dominated by the resolution axis (80.18\% down to 30.16\% across the
quality range) while its salt-and-pepper accuracy never drops below 74.68\%.
It carries no robustness advantage anywhere --- but it is also the fastest
backbone by a factor of $2.5\times$ over DenseNet121 and $6.5\times$ over
TransNeXt-tiny, which matters at deployment.

\sectheading{DenseNet121 --- feature reuse buys little here.} DenseNet121
starts lower on CIFAR-10 (93.56\%) and tracks ResNet50 closely all the way
down (19.58\% at L5), with a small edge at the moderate level (38.54 vs
36.31 at L3, about $2\sigma$ of seed noise). The preliminary hypothesis that
dense feature reuse would help under information loss is not supported: its
Phase~C2 axis ordering and magnitudes are essentially ResNet50's. In plain
terms, keeping copies of earlier features does not help when the features
were destroyed before the first layer ever saw them.

\sectheading{TransNeXt-tiny --- the winner, where winning is possible.}
TransNeXt-tiny posts the best clean baseline on CIFAR-10 (97.64\%), the
largest mild-range margins on the combined curve (+10.11 and +10.41~\Mpp{}
over the CNN mean at L1--L2), and is the most robust backbone on every
Phase~C2 axis at matched image quality --- its blur advantage at L5 is
about 6--12~\Mpp{} over the CNNs. The advantage has a boundary: it shrinks
to +2.05~\Mpp{} at L3 (+0.94 over the better CNN, comparable to the seed
noise) and disappears at L4--L5. On MNIST it never has an advantage at all;
the CNNs are already at ceiling wherever the digits are readable.

\subsubsection{Side-by-side reading}
\label{sec:cm:sidebyside}

Put together, the three profiles say something simple: architecture is a
mild-degradation lever. Where the degraded image still contains texture
(L1--L2 on CIFAR-10, every readable MNIST level), the foveal transformer
converts that texture into up to ten extra accuracy points; where the
texture is gone (L3 and beyond on the combined recipe), all three
architectures converge --- into the 19--21\% band at L5 on CIFAR-10 and the
27--28\% band on MNIST --- and no backbone choice moves the number by more
than the seed noise.

The Phase~C2 attribution explains where the accuracy goes: the resolution
axis carries roughly four times the cost of blur or salt-and-pepper at
matched measured quality. The deployment implication, developed in
Section~\ref{sec:conclusions}, follows directly: pick the transformer when
the operating point is mild and the latency budget allows it, but expect no
classifier --- present or future --- to fix the extreme operating points.
Those need the resolution restored in front of the classifier, not a bigger
classifier behind the sensor.
